\documentclass[oneside,a4paper,14pt]{extarticle}
\usepackage[a4paper,top=20mm,bottom=20mm,left=30mm,right=10mm]{geometry}
\usepackage{amsmath, amssymb, amsfonts, amsthm}
\usepackage[T2A]{fontenc}
\usepackage[utf8]{inputenc}
\usepackage[russian]{babel}
\usepackage{textcomp}
\usepackage{indentfirst}
\usepackage{graphicx}
\usepackage{float}
\usepackage{wrapfig}
\usepackage{caption}
\usepackage{tocloft}
\renewcommand{\cftsecleader}{\cftdotfill{\cftdotsep}}
\usepackage[breaklinks]{hyperref}
\usepackage{titlesec}
\usepackage{fancyvrb}
\usepackage{xcolor}
\titleformat{\section}{\normalsize\bfseries}{\thesection}{1em}{}
\titleformat{\subsection}{\normalsize\bfseries}{\thesubsection}{1em}{}
\titleformat{\subsubsection}{\normalsize\bfseries}{\thesubsubsection}{1em}{}
\renewcommand\baselinestretch{1.33}\normalsize
\setlength{\parindent}{1.25cm}
\emergencystretch=3em
\sloppy
\begin{document}

\setcounter{page}{2}

\section{Цель работы}

Обосновать необходимость разработки веб-платформы для велосипедистов и организаторов спортивных заездов «43Колеса», а также сформировать функциональные и технические требования к ней на основе анализа предметной области и детального сравнения существующих аналогов.

\section{Задачи}

Для достижения поставленной цели необходимо решить следующие задачи:
\begin{itemize}
	\item определить основные понятия предметной области любительского и профессионального велоспорта, обосновать актуальность создания специализированного веб-ресурса;
	\item выделить ключевые потребности и бизнес-процессы пользователей (спортсменов и организаторов), подлежащие автоматизации;
	\item определить источники данных (геолокация, картографические треки, прогнозы погоды, интеграции с социальными сетями);
	\item провести сравнительный анализ существующих решений (сообществ в социальных сетях, афиш, трекеров и канбан-систем), выделив их преимущества и системные недостатки;
	\item сформулировать функциональные и нефункциональные требования к проектируемой веб-системе;
\end{itemize}

\newpage
\section{Анализ предметной области}
\subsection{Предметная область и основные понятия}

Предметной областью проекта является организация, информационное сопровождение и учет участия спортсменов в велособытиях (соревнованиях, любительских заездах, велоквестах и тренировках) на региональном и межрегиональном уровнях (Кировская область и смежные регионы).

Основные понятия предметной области:
\begin{itemize}
	\item \textbf{Веломероприятие (событие)} — запланированный заезд определенной дисциплины (шоссе, кросс-кантри, гревел и т.д.), обладающий регламентом, датой, стартовой локацией, маршрутом и списком участников;
	\item \textbf{Карточка заезда (трек маршрута)} — цифровой профиль маршрута с указанием дистанции, набора высоты, типа покрытия и контрольных точек;
	\item \textbf{временной доска мероприятий} — хронологическое отображение предстоящих и прошедших событий с возможностью фильтрации по дисциплинам и уровню подготовки;
	\item \textbf{Профиль спортсмена} — личная карточка участника, агрегирующая историю преодоленных заездов, пройденный километраж, достижения и статус участия в грядущих стартах;
	\item \textbf{Организатор («Коллега»)} — велоклуб, сообщество или отдельный активист, имеющий право публиковать анонсы заездов, модерировать заявки и публиковать результаты.
\end{itemize}

Объектом автоматизации выступает процесс сбора информации о заездах, регистрации участников и учета их результатов. Предметом автоматизации выступают архитектурные, программные и интерфейсные средства реализации централизованного веб-портала.

\subsection{Актуальность разработки}

С ростом популярности велоспорта и здорового образа жизни количество любительских велосообществ в регионах стабильно увеличивается. Однако информационная инфраструктура данного движения остается фрагментированной:
\begin{itemize}
	\item \textbf{Хаотичность каналов оповещения:} анонсы стартов публикуются в разрозненных беседах мессенджеров из-за чего сложно следить за событиями и планировать участие;
	\item \textbf{Сложность оценки условий старта:} для принятия решения об участии спортсмену требуется вручную сверять несколько независимых источников — текстовое описание заезда, карту с маршрутом (сторонний сервис) и прогноз погоды на дату проведения;
	\item \textbf{Отсутствие сводной истории:} у спортсменов-любителей отсутствует единое портфолио заездов, а учет результатов часто ведется организаторами в файлах Excel.
\end{itemize}

Разработка веб-портала «43Колеса» призвана решить эти проблемы, объединив календарь, картографию, погодный информер и профили велосипедистов в рамках единого интерфейса.

\subsection{Процессы предметной области и роли пользователей}

В системе выделяются три ключевые пользовательские роли:
\begin{enumerate}
	\item \textbf{Пользователь (спортсмен):} регистрация и авторизация; просмотр интерактивной временной доски и карты заездов; подача заявки на участие; просмотр погодных условий; ведение истории участий и личной статистики в кабинете; получение уведомлений через бота;
	\item \textbf{Организатор («Коллега»):} создание и редактирование карточек мероприятий; импорт треков маршрутов; управление списками зарегистрированных участников; кросс-постинг анонсов из сообществ;
	\item \textbf{Администратор:} модерация профилей и карточек мероприятий; предотвращение публикации спама и недостоверных данных; управление правами доступа.
\end{enumerate}

\section{Анализ аналогов}

В ходе предпроектного анализа были изучены существующие программные инструменты и информационные ресурсы, применяемые велосипедистами и организаторами событий.

\subsection{Тематические сообщества ВКонтакте}
Крупнейшие региональные группы и паблики являются основным местом коммуникации велосообществ.
\begin{itemize}
	\item \textit{Достоинства:} огромный охват целевой аудитории, возможность живого комментирования, встроенные функции социальной сети.
	\item \textit{Недостатки:} «умная лента» скрывает анонсы; хронология публикаций не подходит для планирования сезона; отсутствие структурированных карточек соревнований, геофильтрации треков и личной спортивной статистики.
\end{itemize}

\subsection{Портал velogorod.online}
Специализированный веб-архив и каталог веломаршрутов и заездов.
\begin{itemize}
	\item \textit{Достоинства:} наличие базы маршрутов и архивных отчетов о поездках, тематическая направленность.
	\item \textit{Недостатки:} статичность подачи информации;отсутствие удобной системы личных кабинетов с динамической статистикой и интерактивной временной доской ближайших мероприятий.
\end{itemize}

\newpage

\subsection{Trello (и аналогичные канбан-доски)}
Используется активистами для внутренней координации подготовки соревнований и публичных списков заездов.
\begin{itemize}
	\item \textit{Достоинства:} наглядное представление карточек, колоночная группировка, гибкие метки и дедлайны.
	\item \textit{Недостатки:} абсолютная неприспособленность для массового конечного пользователя (спортсмена); отсутствие геопривязки, карт, профилей атлетов и погодных интеграций.
\end{itemize}

\subsection{MapMagic}
Профессиональный сервис построения и визуализации треков для туризма и циклических видов спорта.
\begin{itemize}
	\item \textit{Достоинства:} превосходная картографическая основа, поддержка высотных профилей, импорт/экспорт GPX.
	\item \textit{Недостатки:}специализация инструмента - только редактор маршрута; отсутствие социального взаимодействия, системы регистраций на события и календаря заездов.
\end{itemize}

\newpage
\section{Общие недостатки аналогов и требования к разрабатываемой системе}

\subsection{Общие недостатки существующих решений}
\begin{itemize}
	\item \textbf{Фрагментарность пользовательского пути:} спортсмену необходимо одновременно использовать 3--4 несвязанных инструмента (соцсети для чтения анонсов, картографический сервис для просмотра трека, погодные сайты и таблицы со списками стартующих);
	\item \textbf{Сложность долгосрочного планирования:} отсутствие единой интерактивной временной доски мероприятий с оперативным статусом регистрации;
	\item \textbf{Слабая локализация под региональные нужды:} глобальные сервисы либо заблокированы, либо не поддерживают небольшие любительские старты в регионах.
\end{itemize}

\subsection{Требования к проектируемой системе «43Колеса»}
На основе анализа выявленных проблем сформулированы требования к веб-платформе:
\begin{enumerate}
	\item Разработка адаптивного веб-интерфейса на базе реактивного фреймворка Solid.JS с поддержкой мобильных устройств и десктопов;
	\item Реализация интерактивной временной доски мероприятий с быстрым переключением представлений (список, календарь, карточки);
	\item Интеграция картографического модуля для отображения стартовых городков и маршрутов заездов;
	\item Внедрение виджета погоды с автоматическим запросом метеорологического прогноза на дату проведения заезда в заданной географической точке;
	\item Ведение личного кабинета велосипедиста с историей завершенных заездов и сводной статистикой;
	\item Интеграция с ботом социальной сети «ВКонтакте» для дублирования анонсов и отправки напоминаний о стартах;
	\item Надежный бэкенд на Python с реляционной базой данных PostgreSQL, контейнеризацией через Docker и контролем версий в Git.
\end{enumerate}

\section{Вывод}

В ходе выполнения лабораторной работы была всесторонне исследована предметная область цифровых сервисов для велосипедистов и организаторов спортивных мероприятий. Выявлена актуальность разработки специализированного информационного портала «43Колеса», обусловленная высокой разобщенностью существующих каналов информации.

Проведен сравнительный анализ четырех ключевых аналогов. На основе выявленных ограничений составлен перечень функциональных и нефункциональных требований, а также технологический стек.
\end{document}